\documentclass[12pt, a4paper]{article}

% --- Pacotes Básicos ---
\usepackage{array}
\usepackage[utf8]{inputenc}
\usepackage[T1]{fontenc}
\usepackage[main=brazilian, english]{babel}
\usepackage{indentfirst}
\usepackage{geometry}
\geometry{left=3cm, top=3cm, right=2cm, bottom=2cm}
\usepackage{booktabs}

% --- Informações do Documento ---
\title{Universidade Federal da Fronteira Sul\\
Relatório Construção de um Analisador Léxico}
\author{Erickson Giesel Müller}
\date{\today}

\begin{document}
\maketitle
\thispagestyle{empty} % primeira página não tem número.
%\newpage


\section{Introdução}
	O analisador léxico atua como a primeira fase de um compilador. Sua função é ler cada caractere do código bruto, agrupá-los em unidades lógicas com significado, chamadas lexemas, e classificá-los em categorias, chamadas tokens. O algoritmo de reconhecimento gera uma tabela de símbolos, que armazena atributos de cada lexema lido, e uma fita de saída, que representa a ordem de tokens lidos e será usada como entrada para o sintático.

	O presente trabalho tem como objetivo implementar um analisador léxico baseado no autômato finito determinístico (AFD) gerado através de uma biblioteca de tokens e gramáticas regulares. O analisador lê os caracteres em sequência e separa-os em lexemas, após o reconhecimento dos tokens, armazena na tabela de símbolos e acrescenta o identificador do token à fita de saída.

\section{Referencial Teórico}
	A função de um compilador é traduzir o código-fonte, escrito em linguagem humana, para um programa escrito em linguagem de máquina. A análise léxica separa oa palavras através dos caracteres separadores  e avalia se os caracteres digitados formam palavras que pertencem ao alfabeto da linguagem. Caso não reconheça, informa um erro para ser tratado no token.

	Para que o computador consiga reconhecer essas gramáticas, utiliza-se a Teoria dos Autômatos. Inicialmente, as regras formam um AFND, que pode conter ambiguidades e transições vazias ($\varepsilon$). Como algoritmos de reconhecimento exigem previsibilidade, aplica-se a determinização do autômato. No AFD, para cada estado e cada símbolo lido, existe apenas uma transição de destino possível.	

	A análise léxica também é responsável por lidar com entradas inválidas. Caso o algoritmo leia um caractere que não possui transição mapeada no AFD, o autômato é direcionado para um estado de erro. O analisador reporta a falha, insere o rótulo de erro na tabela de símbolos, registra a posição do lexema em que aconteceu o erro léxico e sincroniza a leitura até o próximo delimitador, permitindo que a compilação continue sem interrupções abruptas.
\section{Implementação e Resultados}
	A construção do automato finito segue as regras determinadas a partir do arquivo \textit{fonte.txt}. A exemplo:\\

\begin{center}
	\textbf{Arquivo \textit{fonte.txt}}\\
\begin{tabular}{|>{\raggedright\arraybackslash}p{8.5cm}|} 
\toprule
se\\
sai\\
foi\\
$S::=f<A>|a<A>|e<A>|i<A>$\\
$A::=f<A>|a<A>|e<A>|i<A>|\varepsilon$\\
\bottomrule
\end{tabular}
\\
\end{center}

	Assim, temos o alfabeto $(\Sigma): \{s, e, a, i, f, o\}$ e o seguinte AFD:
\begin{center}
	\textbf{Autômato Finito Determinístico}\\
\begin{tabular}{|>{\raggedright\arraybackslash}p{1cm}|c|c|c|c|c|c|} 
    \toprule 
    & s & e & a & i & f & o \\
    \midrule 
    $\to S $   & $\{AC\}$ & I    & I    & I    & $\{FI\}$ & $\_$ \\
    $\{AC\}$   & $\_$     & B    & D    & $\_$ & $\_$     & $\_$ \\
    $*B$       & $\_$     & $\_$ & $\_$ & $\_$ & $\_$     & $\_$ \\
    $D$        & $\_$     & $\_$ & $\_$ & E    & $\_$     & $\_$ \\
    $*E$       & $\_$     & $\_$ & $\_$ & $\_$ & $\_$     & $\_$ \\
    $*\{FI\}$  & $\_$     & I    & I    & I    & I        & G    \\
    $G$        & $\_$     & $\_$ & $\_$ & H    & $\_$     & $\_$ \\
    $*H$       & $\_$     & $\_$ & $\_$ & $\_$ & $\_$     & $\_$ \\
    $*I$       & $\_$     & I    & I    & I    & I        & $\_$ \\
    $\_$       & $\_$     & $\_$ & $\_$ & $\_$ & $\_$     & $\_$ \\
    \bottomrule 
\end{tabular}
\end{center}

	A partir da função \textit{analisadorLex()}, o algoritmo de reconhecimento faz a leitura do arquivo \textit{entrada.txt}, analisando as transições de acordo com o AFD gerado.
	($\textbackslash n$ e espaço em branco)

\section{Conclusão e Referências}

\end{document}
